\documentclass[11pt,letterpaper]{article}

% ----- Packages -----
\usepackage[margin=1in]{geometry}    % Page margins
\usepackage{enumitem}               % Customizable lists
\usepackage{hyperref}               % Hyperlinks in PDF
\usepackage{setspace}               % Line spacing (optional)
\usepackage{fontawesome5}
\usepackage{xcolor}                 % Text color
\usepackage[utf8]{inputenc}
\usepackage[T1]{fontenc}
\usepackage{titlesec}
% ----- Document Start -----

\usepackage{color}
\usepackage{fancyhdr}
\usepackage{hyperref}
\usepackage{ifthen}

% \usepackage[yyyymmdd]{datetime}
% \renewcommand{\dateseparator}{-}

% Link formatting.
\definecolor{numcolor}{rgb}{0.5,0.5,0.5}
\definecolor{linkcolor}{rgb}{0,0,0.4}
\definecolor{niceblue}{HTML}{538acf}
\hypersetup{%
    colorlinks=true,        % false: boxed links; true: colored links
    linkcolor=linkcolor,    % color of internal links
    citecolor=linkcolor,    % color of links to bibliography
    filecolor=linkcolor,    % color of file links
    urlcolor=linkcolor      % color of external links
}

% Text formatting.
\newcommand{\foreign}[1]{\textit{#1}}
\newcommand{\etal}{\foreign{et~al.}}
\newcommand{\project}[1]{\textsl{#1}}
\definecolor{grey}{rgb}{0.5,0.5,0.5}
\newcommand{\deemph}[1]{\textcolor{grey}{\footnotesize{#1}}}

% literature links--use doi if you can
  \newcommand{\doi}[2]{\emph{\href{http://dx.doi.org/#1}{{#2}}}}
  \newcommand{\ads}[2]{\href{http://adsabs.harvard.edu/abs/#1}{{#2}}}
  \newcommand{\isbn}[1]{{\footnotesize(\textsc{isbn:}{#1})}}
  \newcommand{\arxiv}[1]{{\href{http://arxiv.org/abs/#1}{arXiv:{#1}}}}

% Section headings.
\renewcommand\familydefault{\sfdefault}
\usepackage{titlesec}

\titleformat{\subsection}
{\normalfont\sffamily\large\bfseries}
{}{0pt}{}[\vspace{-15pt} \enspace \leaders\hrule height 0.4pt \hfill\kern0pt]

\titleformat{\subsubsection}
{\normalfont\sffamily\bfseries}
{}{0pt}{}

\titlespacing{\subsection}{0pt}{2\parskip}{0pt}
\titlespacing{\subsubsection}{0pt}{\parskip}{0pt}

\newcommand{\cvheading}[1]{%
  \addvspace{1ex}%
  \pagebreak[2]%
  \par\noindent
  % Print the heading text
  \textbf{#1}\enspace
  % Insert a rule that fills all remaining horizontal space
  \leaders\hrule height 0.4pt \hfill\kern0pt
  \nopagebreak
  \vspace{-0.4em}%
}

% Set up the custom unordered list.
\newcounter{refpubnum}
\newcommand{\cvlist}{%
    \rightmargin=0in
    \leftmargin=0.15in
    \topsep=0ex
    \partopsep=0pt
    \itemsep=0.2ex
    \parsep=0pt
    \itemindent=-1.0\leftmargin
    \listparindent=0.0\leftmargin
    \settowidth{\labelsep}{~}
    \usecounter{refpubnum}
}

% Set up the custom bullet list for descriptions.
\newlist{bulletdescription}{itemize}{2}
\setlist[bulletdescription]{noitemsep, topsep=-5pt, parsep=0pt, partopsep=0pt, leftmargin=0.25in}
\setlist[bulletdescription,1]{label=$\bullet$}
\setlist[bulletdescription,2]{label=$\circ$}

% Margins and spaces.
\raggedright
\setlength{\oddsidemargin}{0in}
\setlength{\topmargin}{0in}
\setlength{\headsep}{0.20in}
\setlength{\headheight}{0.25in}
\setlength{\textheight}{9.1in}
\addtolength{\topmargin}{-\headsep}
\addtolength{\topmargin}{-\headheight}
\setlength{\textwidth}{6.50in}
\setlength{\parindent}{0in}
\setlength{\parskip}{1ex}

% Headings and footings.
\renewcommand{\headrulewidth}{0pt}
\pagestyle{fancy}
\lhead{\deemph{Adam Boesky}}
\chead{\deemph{Curriculum Vitae}}
\rhead{\deemph{\thepage}}
\cfoot{\deemph{Last updated: \today}}

% Journal names.
\newcommand{\aj}{AJ}
\newcommand{\apj}{ApJ}
\newcommand{\pasp}{PASP}
\newcommand{\mnras}{MNRAS}

\newcommand{\pubsdate}{2026-10-15}
\newcommand{\pubsnumber}{4}
\newcommand{\pubscitations}{67}
\newcommand{\pubshindex}{4}

\addtolength{\topmargin}{-0.4in}
\addtolength{\textheight}{0.8in}
\addtolength{\oddsidemargin}{-0.35in}
\addtolength{\evensidemargin}{-0.35in}
\addtolength{\textwidth}{0.7in}

\begin{document}\thispagestyle{empty}\sloppy\sloppypar\raggedbottom

\textbf{\Large Adam Boesky} \\[0.5ex]
Astrophysics PhD Student and NSF Graduate Research Fellow, Harvard University \\
\textsf{\small %
\href{mailto:apboesky@gmail.com}{apboesky@gmail.com} | %
\href{tel:+1-917-428-3632}{(917)\,428-3632} | %
\href{https://adamboesky.github.io}{adamboesky.github.io} | %
\href{https://github.com/Adam-Boesky}{github.com/Adam-Boesky} | %
\href{https://www.linkedin.com/in/adam-boesky}{linkedin.com/in/adam-boesky} %
}\\[0.5ex]

\vspace{1ex}
\begin{list}{}{\cvlist}
    \item As of \pubsdate\footnote{Publication data collected from \href{https://ui.adsabs.harvard.edu/}{NASA Astrophysics Data System}}, published \pubsnumber\ refereed papers (3 first-author) with \pubscitations\ total citations, and an h-index of \pubshindex; a complete listing can be found on \href{https://ui.adsabs.harvard.edu/public-libraries/Ik0OjllsTWCtf7RzDpNJFQ}{ADS}.
    \item NSF Graduate Research Fellow (2026); Finalist for the Hertz Fellowship, Rhodes Scholarship, and Marshall Scholarship (2025).
    \item Lead developer of \href{https://github.com/Adam-Boesky/astro_SPLASH}{\project{SPLASH}} and \href{https://github.com/Adam-Boesky/ztforce}{\project{ztforce}}; contributor to \href{https://github.com/TeamCOMPAS/COMPAS}{\project{COMPAS}} and \href{https://github.com/alexandergagliano/Prost}{\project{Prost}}; a full list can be found on \href{https://github.com/Adam-Boesky}{GitHub}.
\end{list}

\subsection{Experience}

\begin{list}{}{\cvlist}
    \item
        \textbf{Center for Astrophysics | Harvard \& Smithsonian, Cambridge, MA} \\
        {\color{numcolor} Aug 2023 -- present}\,
        \emph{Researcher, Villar Time-Domain Astronomy Data Lab} \\
        {\color{numcolor} Jun 2022 -- present}\,
        \emph{Researcher, Berger Cosmic Transient Group} \\
        Leading the first multi-survey search for decade-long transients: compiled a catalog of candidates and am now characterizing their physical nature.
        Developed, released, and maintain \href{https://github.com/Adam-Boesky/astro_SPLASH}{\project{SPLASH}}, a machine learning pipeline that classifies supernovae from the inferred properties of their host galaxies at $\sim$500 objects per second, built for the alert stream of the Rubin Observatory's LSST (first-author, \emph{ApJS} 2026).
        Ran binary population synthesis simulations with \project{COMPAS} to model the cosmic merger rates of black holes and neutron stars, showing that they deviate by up to $10\times$ from commonly assumed star-formation-based models (two first-author \emph{ApJ} papers, 2024).
        Contributed to \project{Prost}, a probabilistic host-galaxy association code, and \project{RESOLVE}, a surrogate likelihood for gravitational-wave parameter estimation.
        \\[1.5ex]

    \item
        \textbf{SpaceX, Los Angeles, CA} \\
        {\color{numcolor} Summers 2023, 2024}\,
        \emph{Guidance, Navigation, \& Control Engineering Intern} \\
        Developed physical and empirical models for high-fidelity Monte Carlo simulation of the Falcon 9 and Falcon Heavy rockets.
        Drove operational changes to improve vehicle reliability and performance. Automated day-of-launch operations, built tools to accelerate launch cadence, and assessed the value of mission-specific launch analyses.
        \\[1.5ex]

    \item
        \textbf{ProDex Labs, New York, NY} \\
        {\color{numcolor} Mar 2025 -- Mar 2026}\,
        \emph{Founding Engineer} \\
        Developed a discrete-event simulation engine for building digital twins of manufacturing processes.
        Designed and implemented the planning and scheduling optimizers.
        Integrated an AI agent into the application, speeding up user analysis by $>100\times$. Led development of several modules and managed teams of engineers.
        \\[1.5ex]

    \item
        \textbf{Harvard School of Engineering \& Applied Sciences, Cambridge, MA} \\
        {\color{numcolor} Jun 2020 -- Aug 2022}\,
        \emph{Research Assistant, C3 Policing Research Team} \\
        Quantified the impact of the Counter Criminal Continuum (C3) policing method across business, crime, education, health, and housing data; presented results to Massachusetts officials.
\end{list}

\subsection{Education}
\begin{list}{}{\cvlist}
  \item
        PhD in progress (2026 -- present), Astrophysics, Harvard University.
  \item
        BA 2025, Astrophysics and Mathematics, Harvard University, \emph{cum laude}.
\end{list}

\subsection{Talks, Teaching \& Service}
\begin{list}{}{\cvlist}
  \item Talks at Orange County Astronomers (2026), APS April Meeting (2024), and Post-Pax-VIII (2022); poster at AAS 245 (2025).
  \item Teaching assistant for Harvard Physics 143a and 12b; research mentor to three high school students, including a 2026 Regeneron Science Talent Search Finalist; referee for \emph{The Astrophysical Journal}.
\end{list}

\end{document}
